
\section{Introduction}

Childhood and early adulthood are fundamental years for cognitive development (\cite*{cunha2006interpreting}). School peers can affect cognitive achievement during these crucial years\footnote{See \cite*{ammermueller2009peer,arcidiacono2012estimating,booij2017ability,carrell2013natural,duflo2011peer,garlick2018academic,hanushek2003does,hoxby2000peer,imberman2012katrina,lyle2009effects,sacerdote2001peer}.}, but the mechanisms are not fully understood, limiting our ability to design policies that can be effective across contexts. \par 
This paper studies why classroom peers can shape the academic achievement of students. It introduces a new dataset that links students' academic outcomes, self-reported cost of study effort, teacher curriculum coverage, and schools' resource allocation to newly constructed measures of disruptions to students' environment stemming from one of the most violent earthquakes ever recorded. The dataset allows me to study how study disruptions spill over to classmates. Identification relies on variation in peer disruptions rather than in peer characteristics, avoiding confounding effects associated with the latter.\footnote{Examples of studies using naturally occurring exogenous variation in peer characteristics to account for confounding effects are \cite{hoxby2000peer, angrist2004does,hoxby2005taking,lavy2012inside,imberman2012katrina}.} I draw on the empirical evidence to propose a new theory of peer influence in cognitive development. While formulated within the context of an environmental shock, the theory offers a framework to understand why peers matters for learning in many contexts.\par 



The empirical context is the 2010 Maule mega-earthquake, the seventh strongest ever instrumentally recorded \citep{USGSCatalog2025}. I combine administrative and survey data with information on damage propagation among over 150,000 students. I start by building a measure of each student's home's vulnerability to the earthquake using information on housing quality. From the last pre-earthquake census I obtain information on the construction materials of the homes of the nearly one million Chilean households with at least one school-aged child. I employ an unsupervised learning algorithm to stochastically assign their homes to seismic resistance classes. Armed with this housing quality measure, I develop a model that can accurately predict housing quality from a household's characteristics, and apply it to administrative data on Chilean students to predict the quality of their homes. Drawing upon the structural engineering literature, I then combine this newly developed measure of housing quality with geocoded information on ground-shaking intensity in each student's hometown to construct a measure of home damages for each student.\footnote{I gratefully acknowledge Prof. Sergio Ruiz of the Geology Department at the University of Chile, a leading expert on the seismic vulnerability of Chilean buildings, for feedback on the damage measure. The measure falls within Deterministic Earthquake Loss Estimation, which aims to estimate losses from a specific seismic event. It stands in contrast to Probabilistic Earthquake Loss Estimation, which aims to predict potential losses from many possible seismic events (\cite{McGuire2004Seismic}).} This variable measures the shock to each student's environment.\par



I link the damage measure to administrative and survey data from the Chilean Ministry of Education. Data on students include standardized test scores and GPA at two points in time (in fourth and eighth grade), family background, type of school attended, and survey information on self-reported cost of study effort and ability to engage with course content. Data on teachers include the fraction of the curriculum they were able to cover. For the 42\% of schools participating in the preferential school subsidy program (SEP, {\itshape Subvenci\'{o}n Escolar Preferencial}), data include detailed reports on SEP resource spending.  All observations can be assigned to classrooms and schools through unique identifiers.\par  


Using this newly constructed dataset, I first document new facts about socioeconomic segregation among Chilean children. Poorer students were more likely to live in the rural localities experiencing more ground shaking, and conditional on locality, their homes were built according to worse construction standards. As a result, compared with students whose parents have more than 14 years of education (who likely attended some college), those whose parents have at most 14 incurred twice the amount of home damages (USD 1,552 vs. USD 759, or 47\% vs. 23\% of annual household income). \par 

I then estimate the causal impacts on students' outcomes of damages to their own homes and to the homes of their classmates, focusing on the mean and standard deviation of damages among classmates. The identification strategy relies on a difference-in-differences framework leveraging the different correlation between seismic vulnerability and outcomes across two cohorts of students, one whose outcomes were measured in 2009, before the earthquake struck, and one whose outcomes were measured in 2011, after the earthquake struck. This strategy eliminates confounding effects that could arise if the measures of earthquake vulnerability correlated with unobserved outcome determinants. The identifying assumption is that the relationship between outcomes and earthquake vulnerability would be the same across cohorts absent the earthquake. I provide several pieces of evidence supporting this assumption, showing no evidence that students reallocated to classrooms and schools in response to the earthquake in the estimation sample (which excludes by design forced re-locations), showing no impacts on placebo outcomes pre-determined at the time of the earthquake, and showing that the seismic vulnerability of peers' homes does not correlate differently with outcomes across cohorts in regions that were never affected by the earthquake.\par 

Using this strategy, I find that the damage incurred by a student's own home had a negative and non-negligible effect on achievement 22 months post-earthquake. A 1 standard deviation increase in damages lowered test scores by 0.03 standard deviations, and GPA by 0.02 standard deviations, albeit the GPA impacts are statistically insignificant. Using survey data, I provide evidence that damages at the students' own homes increased students' reported cost of study effort, suggesting it could have mediated the detrimental impacts on achievement. \par
Regarding the spillover effect of damages to peers' homes keeping fixed a student's own exposure to the earthquake, I find that increasing the average damages suffered by classroom peers increases own test scores, and the effects are not significantly different across students with different initial performance. This appears to be the result of schools overcompensating any potential negative impacts. Data on school spending show that schools responded to the average level of damages suffered by their students by reallocating resources from recruitment costs toward activities directly linked to student support and learning recovery, such as educational and psychological support.\par 
In contrast, increases in the within-classroom standard deviation of damages lowered the achievement of students with high initial performance and increased the achievement of those with low initial performance. Neither the data on curriculum coverage nor the data on school spending provide statistically significant evidence of schools reacting to how dispersed the damages were, such as by focusing existing resources towards activities targeted at students with lower initial performance. Additionally, emergency funds were granted to schools depending on the overall damage severity, not its dispersion (\cite{Gobierno2010}, Appendix \ref{sec:reconstructionplan}). \par 
The results so far suggest that study disruptions had spillover effects on classroom peers, and that schools mediated the effects of average damages but possibly not of damage dispersion, suggesting a potential role for peer-to-peer interactions. To better understand the damage-dispersion spillovers, I analyze impacts on GPA rank, to examine if the heterogeneous effects on GPA triggered changes to students' GPA rankings in the classroom. Surprisingly, I find this did not occur. Students with higher initial performance experienced drops in GPA in classrooms with more dispersed damages, without an accompanying drop in their GPA rank.  A possible reason for this is that students care about their GPA rank. Faced with a changed cost of study effort among their peers, students adjusted their effort and learning (a peer effect), but not at the expense of their classroom standing in terms of an achievement measure observable to classmates. Rank concerns, therefore, could be a mode of peer-to-peer interactions.\par

 
 

Drawing on this insight and on the empirical findings, I formulate a new theory of interactions in the classroom. I formalize the simple intuition that students care about their classroom standing through a game-of-status model of simultaneous effort choices in the classroom. In the model, students are characterized by an effort-cost type, which is affected by prior test scores, socioeconomic characteristics, and by the damages to their home. They produce GPA by exerting costly effort, and derive utility from GPA and GPA rank. To account for the evidence on schools' reactions, mitigating actions by schools in response to average damages are allowed to enter the achievement production technology. I show that the model, which admits a unique symmetric Bayesian Nash equilibrium, can rationalize the empirical findings. Specifically, a school's mitigating efforts lead to positive effects of average damages. The competition motive behind students' effort choices causes the damage dispersion to have heterogeneous impacts on GPA, and null effects on GPA rank, along the prior test score distribution, conditional on a student's own damage and socioeconomic characteristics. By changing the density of types differently across the effort–cost type distribution, damage dispersion affects incentives differently for different students, both by changing how many nearby types can be overtaken with a marginal increase in effort, and by shifting equilibrium effort responses throughout the distribution. These forces generate heterogeneous effects on the returns to effort and, in turn, on GPA, that are consistent with the patterns found in the data.\footnote{A small strand of the literature on college admissions has developed tournament models of student effort under rank incentives (\cite{bodoh2017pre,grau2018impact,tincani2025college}). These papers study the effects of {\itshape changing} the rank incentives, holding peer characteristics fixed, rather than peer effects {\itshape given} rank concerns. The most relevant is \cite{tincani2025college}: after showing experimentally that rank incentives affect study effort in Chilean high schools, the authors develop and structurally estimate a tournament model of simultaneous effort choices, in which college seats are assigned according to within-school GPA rank.\label{fn:tournaments} }  \par 

The central theoretical insight can be applied more broadly outside the realm of environmental shocks: when competitive motives drive study effort, changing the dispersion of peers' cost-of-effort types, be it through shocks to the cost of effort or through compositional changes from classroom assignment policies, affects learning, and does so differently for different students depending on how the change affects the number of nearby competitors and the effort of all competitors. This has important and so far mostly unexplored implications for policy, as I discuss in this article's concluding section.\par 


Methodologically, this study relates to the small literature that examines peer interactions relying on random shocks to students that keep group composition constant (see the survey in \cite{bramoulle2020peer}). One of the closest studies is \cite{fruehwirth2013identifying}, who exploits the introduction of a student accountability policy in North Carolina targeted at low-achievers. The policy serves as a shock to the effort of some but not all students in the classroom; the fraction of affected peers is used to estimate the impact of peers' achievement on own achievement within a linear-in-means framework.\footnote{\cite{berlinski2023helping} have applied a similar strategy to data from a literacy remediation program in Colombia, and \cite{dieye2014accounting} to data from a randomized experiment on a scholarship program in Colombia. See also \cite{fruehwirth2014can} for an in-depth analysis of the identification of the effect of contemporaneous peer outcomes on own outcomes when outcomes are partly determined by unobserved factors.  } The estimates are interpreted as best-response functions through the lens of a model of effort choices in a strategic environment where students desire to conform to each other. In contrast, this paper considers a continuous shock, the extent of damages that each student's home incurred. Rather than identifying best-response functions, an `endogenous peer effect' in the terminology of \cite{manski1993identification}, this paper examines the reduced-form impact of changing the distribution of the shocks within the classroom, an `exogenous peer effect' in this terminology. Through the lens of a model where students desire to compete with each other, it interprets the estimates as comparative statics on equilibrium outcomes in the classroom.\footnote{Unlike \cite{fruehwirth2013identifying}, I do not estimate the `endogenous peer effect', because the theoretical model does not yield a closed-form expression for the best-response function. Instead, the model characterizes the equilibrium effort function within each classroom, which traces effort as a function of a student's effort-cost type, and derives comparative statics on this equilibrium effort function under different within-classroom distributions for the effort-cost types. When damages to a student's home affect the student's effort-cost type, the model provides testable implications on the shape of the spillovers from earthquake-induced disruptions.} The impacts of the mean and of the dispersion of this continuous shock among peers are shown to be helpful to inform a new theory of peer influence.\footnote{\cite{de2013understanding} also test theories of peer influence. They use evidence on the effects of changing peer composition, which is deliberately kept constant here.}\par 

Within the vast literature on peer effects in education, relatively few studies have developed theories of peer influence. Existing theories commonly assume that students have a desire to conform to their peers, or that there are complementarities between peers in the achievement production technology (e.g. \cite*{brock2001discrete,durlauf2006multinomial,calvo2009peer,fruehwirth2013identifying,conley2024social}). Both assumptions rationalize the workhorse linear-in-means model of peer effects with continuous outcomes (\cite*{blume2015linear}). In contrast, I present a new theory that rationalizes why moments beyond the mean may matter, in a parsimonious way. It offers a simple insight: when students derive utility from rank, changing the cost of study effort of peers affects own effort, because it changes the ability of peers to compete. Empirically, this generates a peer effect where moments beyond the mean matter, without the need to introduce extensions to the technology or preferences to capture the influence of higher-order moments. Such a mechanism has been largely ignored despite its intuitive appeal.\footnote{In line with the theory first introduced in this paper (as detailed in e.g. \cite{tincani2017heterogeneous}), \cite{rosenzweig2023classroom} recently provided evidence supporting this mode of interaction within the context of Southeast Asian refugee students in the US.} \par 
 
The idea that students may care about their rank is consistent with a growing body of evidence showing that competitive preferences can emerge early in life (\cite{sutter2015gender,page2017older}) and that classroom rank can yield both immediate and longer-term benefits. A large literature on {\itshape rank effects} has examined how exogenous changes in a student’s achievement rank within a reference group affect student outcomes. A higher rank has been shown to improve self-concept, happiness, and teacher perceptions of ability, and to raise later educational attainment and earnings (e.g. \cite{zeidner1999big,marsh2007big,elsner2017big,elsner2018rank,murphy2020top,denning2023class,ladant2024essential,carneiro2025effect}). Students, therefore, may value their relative standing, even when it is not formally rewarded. Building on this insight, this paper studies instead the implications of {\itshape rank concerns}, i.e., of students caring about their rank, for how peers influence each other's learning.\par  


This paper also relates to the empirical literature using natural disasters to identify peer effects in education, such as \cite{cipollone2007social}, \cite{sacerdote2008saints}, and \cite{ imberman2012katrina}. In contrast to these studies, it does not use forced relocations of students for identification, focusing instead on variation in the intensity of disruptions.\footnote{This distinguishes this paper also from the experimental and quasi-experimental literature that use variation in assignment to peer groups, such as dorms (\cite{sacerdote2001peer,zimmerman2003peer,Stinebrickner2006,kremer2008peer,garlick2018academic}) or classrooms (\cite{duflo2011peer,whitmore2005resource,kang2007classroom}).} I interpret these disruptions as shocks to a primitive of the students' skill-accumulation problem that shape their effort choices. This interpretation, consistent with the skill-formation framework where human capital reflects optimally chosen investments (\cite{todd2003specification,cunha2008formulating,cunha2010estimating}), allows me to move beyond measuring earthquake spillovers to shed light on how social interactions influence those investment decisions. Finally, the paper relates to the empirical literature on ability peer effects studying the impacts of moments beyond the mean (\cite{lyle2009effects,booij2017ability,ding2007peers,vigdor2007peer, hoxby2005taking}) and of partitioning the support of ability, which varies first and higher-order moments simultaneously (\cite{carrell2013natural,duflo2011peer}). These studies tend to find that moments beyond the mean matter for learning. \par 




The article is structured as follows. Section \ref{sec:data_and_measurements} details the data, damage measure, and describes damage propagation among students, documenting new facts about socioeconomic inequalities. Section  \ref{sec:analysis_earthquake_effects} delves into the empirical analysis of damage effects on achievement, and assesses the identifying assumption and robustness. Section \ref{sec:mechanisms} presents evidence on mediating factors using administrative and survey data. Section \ref{sec:new_theory_peer_effects} introduces the theory of peer influence based on rank concerns, rationalizing the evidence. Section \ref{sec:conclusions} concludes, discussing policy implications and suggesting future research avenues. \par 
